\documentclass[10pt,a4paper]{amsart}
\usepackage[margin=19mm]{geometry}
\usepackage{amsmath,amssymb,mathtools}
\usepackage[scr=rsfs,cal=euler]{mathalfa}
\usepackage{xfrac}
\usepackage[unicode,hidelinks,bookmarks=false]{hyperref}
\newcommand{\ie}{i.e.\ }
\newcommand{\cf}{cf.\ }
\newcommand{\resp}{respectively\ }
\newcommand{\Subsection}{\subsection}
\providecommand{\nobreakdash}{\nobreak\hbox{-}}
\setlength{\emergencystretch}{2em}
\multlinegap0pt
\usepackage{xcolor}
\usepackage{amssymb}
\usepackage[shortlabels]{enumitem}
\usepackage{mathtools}
\usepackage{tikz}
\usepackage{bm}
\usepackage{float}
\newtheorem{corollary}{Corollary}
\newtheorem{lemma}{Lemma}
\newtheorem{theorem}{Theorem}
\title{$3$-designs from symmetric conference matrices}
\author{Gary Greaves, Sho Suda}
\date{Source: arXiv:2402.17528v2 (CC BY 4.0)}
\begin{document}
\maketitle

\noindent\emph{Source and licence.} $3$-designs from symmetric conference matrices. Version arXiv:2402.17528v2 (CC BY 4.0), distributed by arXiv under the Creative Commons Attribution 4.0 International licence, http://creativecommons.org/licenses/by/4.0/, as recorded in the arXiv licence metadata for this exact version and shown on the abstract page https://arxiv.org/abs/2402.17528v2. Full licence notice: \texttt{licenses/arxiv-2402.17528-CC-BY-4.0.txt}.\par\smallskip

\noindent\textit{Editorial scope.} Counted unit: the article's theorem thm:pbibd2 (source0.tex lines 1691--1709). The article's complete original proof of it is delivered with it (source0.tex lines 1710--1766, one \texttt{proof} environment of 57 lines). No mathematics is written, repaired or re-derived: the statement and the proof are the article's own lines, in the article's own notation, and no retained line is substituted or re-flowed.  Retained uncounted as the source-local context the proof needs: lines 400--411; lines 426--441.  Nothing was omitted from any retained span, so the package carries no omission record.  No retained line contains a citation, so the wrapper carries no bibliography and no bibliography entry.  No retained line contains a figure environment, an image-inclusion command or drawing code, so no image is delivered and the package declares no asset dependency.  Editorial formatting only, in addition: an AMS \texttt{amsart} preamble that restates the article's own declarations and macros used by the retained lines, namely the environments \texttt{corollary}, \texttt{lemma}, \texttt{theorem} on separate counters, together with the standard packages those lines need.  The retained environments print in this document as Corollary 1, Lemma 1, Theorem 1 in order of appearance, whereas the article prints the same items with the numbering of its own full document.  The complete native source of the article as distributed in the arXiv e-print for this exact version is bundled unchanged as \texttt{sources/155211/source0.tex}.
\begin{corollary}\label{cor:mu}
Let $v$ be a positive integer, $k \in \{0,1,\dots,v\}$, $A$ be a $v\times v$ complex matrix, and let $\alpha$ be a subset of $[v]$. 
Then
\begin{enumerate}
    \item $\displaystyle \binom{v-|\alpha|}{k} = \sum_{d\in D_A(k)}\mu_A(\alpha,k,d)$;
    \item  
%\begin{align*}
     $\displaystyle c_A(\alpha,k)    =(-1)^{k}\sum_{d\in D_A(k)}d\mu_A(\alpha,k,d).$
%\end{align*}

\end{enumerate}
\end{corollary}

\begin{lemma}
    \label{lem:lambda}
        Let $v$ and $t$ be positive integers and $\mathfrak B$ be a family of subsets of $[v]$. 
    For $\beta\subset [v]$, we have
    \begin{align*}
    % \label{eq:lambda}
%    \lambda_j(\beta)=\sum_{h=0}^j (-1)^h \binom{j}{h}\lambda_0^h(\beta).
\lambda_\mathfrak B(\beta)=\sum_{\gamma \subset \beta}(-1)^{|\gamma|} \mu_\mathfrak B(\gamma).
\end{align*}
% If $\lambda_0^j=\lambda_0^j(\beta)$ depends only on $j=|\beta|$ for $1\leqslant j\leqslant t$, then so does 
Suppose that, for each $i \in \{0,1,\dots,t\}$, there exists $\mu_i$ such that $\mu_i = \mu_\mathfrak B(\beta)$ for all $\beta \in \binom{[v]}{i}$. 
Then, for each $i \in [t]$ and $\beta \in \binom{[v]}{i}$, we have
\begin{align*}
    \lambda_\mathfrak B(\beta)=\sum_{j=0}^i (-1)^j \binom{i}{j}\mu_j.
\end{align*}
\end{lemma}

\begin{theorem}\label{thm:pbibd2}
Let $\mathcal{X}=([v],\{R_i\}_{i=0}^d)$ be a symmetric association scheme. 
Let $k$ be a positive integer, and $A$ be a $v\times v$ complex diagonalisable matrix such that  
\begin{enumerate}
    \item $D_A(k)=\{a,b,c\}$, where $|D_A(k)| = 3$,
    \item $C_A \left (\binom{[v]}{i},k \right) = \{c_i\}$ for each $i \in \{0,1\}$ and $C_A(\tilde{R}_j,k) = \{c_2(j)\}$ for each $j \in [d]$, and 
    %\item $\left ([v],\binom{[v]}{v-k}\setminus\mathfrak B_A(v,k,c)\right )$ 
    %is a $\operatorname{PBIBD}(v,v-k;\eta_2(1),\dots,\eta_2(d))$ on $\mathcal{X}$. 
    \item %$\left ([v],\mathfrak B_A(v,k,c)\right )$ 
    the hypergraph $\mathfrak H_{A}(v,k,c)$  is a $\operatorname{PBIBD}(v,k;\eta_2(1),\dots,\eta_2(d))$ on $\mathcal{X}$.
\end{enumerate}
Then the hypergraph $\mathfrak H_{A}(v,k,a)$ 
is a $\operatorname{PBIBD}(v,k;\lambda_1,\dots,\lambda_d)$ on $\mathcal{X}$ where, for each $j \in [d]$, we have
\[
\lambda_j=\frac{(-1)^k}{a-b}\left(c_0-2c_1+c_{2}(j)\right)-\frac{b}{a-b}\binom{v-2}{k-2}-\frac{c-b}{a-b}\eta_2(j).
\]
%where $\eta_{0} = \binom{v}{v-k} -|\mathfrak B_A(v,k,c)|$ and $\eta_1 = \left |\left \{ \alpha \in \binom{[v]}{v-k}\setminus\mathfrak B_A(v,k,c) \; : \; 1 \in \alpha \right \} \right |$.
% where $\eta_{0} = |\mathfrak B_A(v,k,c)|$ and $\eta_1 = \left |\left \{ \alpha \in \mathfrak B_A(v,k,c) \; : \; 1 \in \alpha \right \} \right |$.
\end{theorem}

\begin{proof}
% Fix $x\in\{a,b\}$ and let $y$ be such that $\{x,y\}=\{a,b\}$. 
Let $\mathfrak B = \mathfrak B_A(v,k,a)$, $\eta_{0} = |\mathfrak B_A(v,k,c)|$, and $\eta_1 = \left |\left \{ \alpha \in \mathfrak B_A(v,k,c) \; : \; 1 \in \alpha \right \} \right |$.
% Let $\eta_{0}$ be the number of blocks in the PBIBD $\mathfrak H_{A}(v,k,c)$ and $\eta_1$ be the number of blocks containing a point in the PBIBD $\mathfrak H_{A}(v,k,c)$. 
First we claim that, for each $i\in\{0,1\}$, there exists $\mu_i$ such that $\mu_\mathfrak B(\beta) = \mu_i$ for each $\beta \in \binom{[v]}{i}$. 
% \textcolor{red}{By interchanging the roles of $\lambda$ and $\mu$ in Lemma~\ref{lem:lambda}, one can show that the number of blocks in $\mathfrak H_{A}(v,k,c)$ that do not contain a point is $\eta'_{1}:=v-k$ and that of blocks that do not contain $\{x,y\}\in \tilde{R}_j$ is $\eta'_2(j):=v-2k+\eta_2(j)$.}  
For $x \in [v]$, denote by $\nu_1$ the number of hyperedges of $\mathfrak H_{A}(v,k,c)$ that do not contain $x$.
For $\{x,y\}\in \tilde{R}_i$, denote by $\nu_2(i)$ the number of hyperedges of $\mathfrak H_{A}(v,k,c)$ that do not contain $\{x,y\}$.
Using a straightforward inclusion-exclusion argument, we can write $\nu_1 = v - \eta_1$ and $\nu_2(i) = v - 2\eta_1+\eta_2(i)$ for each $i \in [d]$.
Also, define $\nu_0 = \eta_0$.
By Corollary~\ref{cor:mu}
% (1), (2) 
together with assumptions (1) and (3), for each $i\in\{0,1\}$, we obtain
\begin{align*}
    c_i &= c_A(\beta,k)=(-1)^{k}\left(a\mu_\mathfrak B(\beta)+b\left(\binom{v-i}{k}-\nu_i-\mu_\mathfrak B(\beta)\right)+c\nu_i\right)\nonumber\\
    &= (-1)^{k}\left((a-b)\mu_\mathfrak B(\beta)+b\binom{v-i}{k}+(c-b)\nu_i\right)\text{ for }\beta\in\binom{[v]}{i}.
\end{align*}
from which our claim follows.
Next, for any $\beta \in \binom{[v]}{2}$, there exists $j \in \{1,\dots,d\}$ such that $\beta \in \tilde{R}_j$.
Furthermore, we claim that, for each $j \in \{1,\dots,d\}$, there exists $\mu_2(j)$ such that $\mu_\mathfrak B(\beta) =\mu_2(j)$ for each $\beta \in \tilde{R}_j$.
Again, by Corollary~\ref{cor:mu}
together with assumption (1) and (3), for each $j \in \{1,\dots,d\}$, we obtain
\begin{align*}
    c_2(j) &=  c_A(\beta,k)=(-1)^{k}\left(a\mu_\mathfrak B(\beta)+b\left(\binom{v-2}{k}-\nu_2(j)-\mu_\mathfrak B(\beta)\right)+c\nu_2(j)\right)\nonumber    \\
    &= (-1)^{k}\left((a-b)\mu_\mathfrak B(\beta)+b\binom{v-2}{k}+(c-b)\nu_2(j)\right)\text{ for }\beta\in\tilde{R}_j,
\end{align*}
from which the claim follows.
Hence, for each $i\in\{0,1\}$, we can write $\mu_i=\mu_\mathfrak B(\beta)$ for some $\beta \in \binom{[v]}{i}$ and $\mu_{2}(j)=\mu_\mathfrak B(\beta)$ for some $\beta \in \tilde{R}_j$.

By Lemma~\ref{lem:lambda},
$$
\lambda_\mathfrak B(\beta)=\sum_{\gamma\subset \beta}(-1)^{|\gamma|}\mu_\mathfrak B(\gamma)=\begin{cases}\mu_0 & \text{ if }|\beta|=0,\\
\mu_0-\mu_1 & \text{ if }|\beta|=1,\\
\mu_0-2\mu_1+\mu_{2}(j) & \text{ if }|\beta|=2\text{ and }\beta\in \tilde{R}_j.
\end{cases} 
$$
Thus, $\lambda_\mathfrak B(\beta)$ does not depend on $\beta$ for $\beta \in \binom{[v]}{i}$ when $i\in\{0,1\}$ and depends only on $j$  if $|\beta|=2$ and $\beta\in \tilde{R}_j$.  
Hence, the hypergraph $\mathfrak H_{A}(v,k,a)$ 
is a $\operatorname{PBIBD}(v,k;\lambda_1,\ldots,\lambda_d)$ on $\mathcal{X}$.

Finally, we establish the expressions for the parameters $\lambda_1,\ldots,\lambda_d$.  
Since 
\begin{align*}
\mu_i&=\frac{1}{a-b}\left((-1)^{k} c_{i}-b\binom{v-i}{k}-(c-b)\nu_{i}\right)\quad\text{for $i\in \{0,1\}$},\\  
\mu_{2}(j)&=\frac{1}{a-b}\left((-1)^{k} c_{2}(j)-b\binom{v-2}{k}-(c-b)\nu_{2}(j)\right)\quad\text{for $j\in \{1,\dots,d\}$},   
\end{align*}
the parameters $\lambda_j$ are determined as 
\begin{align*}
\lambda_j&=\mu_0-2\mu_1+\mu_{2}(j)\\
    %&=\frac{1}{a-b}\left((-1)^{k} c_{0}-b\binom{v}{k}\right)-\frac{2}{a-b}\left((-1)^{k} c_{1}-b\binom{v-1}{k}\right)+\frac{1}{a-b}\left((-1)^k c_{2}(i)-b\binom{v-2}{k}\right)\\
    %&=\frac{(-1)^k}{a-b}\left(c_0-2c_1+c_{2}(i)\right)+\frac{b}{a-b}\left(-\binom{v}{k}+2\binom{v-1}{k}-\binom{v-2}{k}\right)\\
    &=\frac{(-1)^k}{a-b}\left(c_0-2c_1+c_{2}(j)\right)-\frac{b}{a-b}\left( \binom{v}{k} - 2\binom{v-1}{k}+\binom{v-2}{k}\right)-\frac{c-b}{a-b}(\nu_{0}-2\nu_1+\nu_2(j))\\
    % &=\frac{(-1)^k}{a-b}\left(c_0-2c_1+c_{2}(j)\right)-\frac{b}{a-b}\binom{v-2}{k-2}-\frac{1}{a-b}(v-2(v-\eta_1)+(v-2\eta_1+\eta_2(j)))\\
    &=\frac{(-1)^k}{a-b}\left(c_0-2c_1+c_{2}(j)\right)-\frac{b}{a-b}\binom{v-2}{k-2}-\frac{c-b}{a-b}\eta_2(j).\qedhere
\end{align*}
% $\eta$ to $\nu$ ****
\end{proof}

\end{document}
